% Part: sets-functions-relations
% Chapter: relations
% Section: reflections
%
\documentclass[../../../include/open-logic-section]{subfiles}

\begin{document}

\olfileid{sfr}{rel}{ref}
\olsection{તત્વજ્ઞાનવિષયક વિચારણા}

\olref[set]{sec}માં આપણે સંબંધોને ચોક્કસ ગણો તરીકે વ્યાખ્યાયિત
કર્યા હતા. અહીં થોડું થોભીને એક તત્વજ્ઞાનવિષયક પ્રશ્ન પૂછવો જોઈએ:
આવી વ્યાખ્યા શું \emph{કરે છે}? આપણે કોઈ તત્વમીમાંસાકીય
તાદાત્મ્યનાં તથ્યો \emph{શોધી કાઢ્યાં} છે એમ કહેવું જોઈએ,
તે ભારે શંકાસ્પદ છે; દાખલા તરીકે, $\Nat$ પરનો ક્રમ સંબંધ
આપણે \olref[set]{sec}માં વ્યાખ્યાયિત કરેલો ગણ
$R=  \Setabs{\tuple{n,m}}{n, m \in \Nat\text{ અને }
n < m}$ \emph{જ નીકળ્યો} એમ કહેવું. તેના ત્રણ કારણો છે.

પ્રથમ: \olref[set][pai]{wienerkuratowski}માં આપણે
$\tuple{a, b} = \{\{a\}, \{a, b\}\}$ વ્યાખ્યાયિત કર્યું હતું.
તેને બદલે $\lVert a, b\rVert = \{\{b\}, \{a, b\}\} = \tuple{b,a}$
એવી વ્યાખ્યા લો. જ્યારે $a \neq b$, ત્યારે
$\tuple{a, b} \neq \lVert a,b\rVert$ થાય છે. પણ
$\tuple{a,b}$ને બદલે $\lVert a,b\rVert$ને ક્રમયુક્ત જોડની
આપણી વ્યાખ્યા તરીકે એટલી જ સારી રીતે સ્વીકારી શક્યા હોત.
બંને વ્યાખ્યાઓ સરખી રીતે કામ આપત. તેથી હવે આપણી પાસે
પ્રાકૃતિક સંખ્યાઓ પરનો ક્રમ સંબંધ “હોવા” માટે બે સરખા
યોગ્ય ઉમેદવારો છે, એટલે કે:
\begin{align*}
		R &= \Setabs{\tuple{n,m}}{n, m \in \Nat \text{ અને }n < m}\\
		S &= \Setabs{\lVert n,m\rVert}{n, m \in \Nat \text{ અને }n < m}.
\end{align*}
ઘટકો દ્વારા નિર્ધારિત સમાનતાના સિદ્ધાંતથી $R \neq S$ હોવાથી,
સ્પષ્ટ છે કે તે \emph{બંને} $\Nat$ પરના ક્રમ સંબંધ સાથે
તાદાત્મ્ય ધરાવી શકે નહીં. પરંતુ હકીકતમાં $S$ને બદલે $R$
(અથવા ઊલટું) \emph{જ} ક્રમ સંબંધ છે એવો દાવો કરવો
મનસ્વી અને તેથી કંઈક ક્ષોભજનક બને. (આ ગણસિદ્ધાંતમાં બધું
ઘટાડી શકાય એવા મત વિરુદ્ધની દલીલનો એક અત્યંત સરળ દાખલો છે,
જેને બેનાસેરાફે \citeyear{Benacerraf1965}માં જાણીતી બનાવી હતી.
આપણે ઘણી વાર ફરી તેની ચર્ચા કરીશું.)

બીજું: જો આપણે માનીએ કે \emph{દરેક} સંબંધને કોઈ ગણ સાથે
તાદાત્મ્ય ધરાવતો ગણવો જોઈએ, તો ગણના સભ્યપદનો સંબંધ $\in$
પોતે એક ચોક્કસ ગણ હોવો જોઈએ. ખરેખર, તે
$\Setabs{\tuple{x,y}}{x \in y}$ ગણ જ હોવો જોઈએ.
પરંતુ શું આ ગણ અસ્તિત્વ ધરાવે છે? રસેલના વિરોધાભાસને જોતાં,
આવો ગણ અસ્તિત્વ ધરાવે છે એવો દાવો સહજ રીતે સ્વીકારી શકાય તેમ નથી.
ખરેખર, \oliflabeldef{cumul:::part}{આપણે
\olref[cumul][][]{part}માં વિકસાવવાનો ગણસિદ્ધાંત આ ગણના
અસ્તિત્વનો \emph{નિષેધ} કરશે.\footnote{આગળની વાત અહીં
કરીએ તો, કારણ આ છે. પરસ્પરવિરોધ મેળવવા માટે ધારો કે
$I = \Setabs{\tuple{x,y}}{x \in y}$ અસ્તિત્વ ધરાવે છે.
તો $\bigcup \bigcup I$ સાર્વત્રિક ગણ છે, જે
\olref[sfr][z][sep]{thm:NoUniversalSet}નો વિરોધ કરે છે.}}{ગણસિદ્ધાંતને
સ્વયંસિદ્ધિઓ પર આધારિત સિદ્ધાંત તરીકે ચોકસાઈપૂર્વક વિકસાવી
શકાય છે, અને તે સિદ્ધાંત ખરેખર આ ગણના અસ્તિત્વનો નિષેધ કરશે.}
તેથી, કેટલાક સંબંધોને ગણ તરીકે લઈ શકાય તો પણ, ગણના
સભ્યપદના સંબંધને વિશેષ કિસ્સો ગણવો પડશે.

ત્રીજું: સંબંધોને ગણો સાથે “એકરૂપ” ગણતી વખતે આપણે કહ્યું હતું
કે $\tuple{x,y} \in R$ માટે $Rxy$ લખવાની છૂટ રાખીશું.
આ બરાબર છે, જો સભ્યપદના સંબંધ “$\in$”ને વિધેય
\emph{તરીકે} લેવામાં આવે. પરંતુ જો આપણે માનીએ કે “$\in$”
કોઈ ચોક્કસ પ્રકારના ગણને સૂચવે છે, તો “$\tuple{x,y} \in R$”
એ પદપ્રયોગમાં ગણોને સૂચવતાં માત્ર ત્રણ એકવચની પદો જ છે:
“$\tuple{x,y}$”, “$\in$” અને “$R$”.
આવી નામોની યાદી કોઈ વિધાન વ્યક્ત કરવા માટે
“પ્યાલો કલમદાની મેજ” જેવી નિરર્થક પ્રતીકશ્રેણી કરતાં વધુ
સમર્થ નથી. ફરીથી, કેટલાક સંબંધોને ગણ તરીકે લઈ શકાય તો પણ,
ગણના સભ્યપદનો સંબંધ વિશેષ કિસ્સો હોવો જોઈએ.
(આમાં ફ્રેગેના \emph{ઘોડો} ખ્યાલના વિરોધાભાસનું સરળ સ્વરૂપ
અને વિટ્ગેન્સ્ટાઇને એક વખત રસેલ સામે ઉઠાવેલો પ્રસિદ્ધ વાંધો
ભેગાં થયાં છે.)

તો આ પરથી આપણે શું સમજવું? સંખ્યાઓ પરના સંબંધો ગણો છે એમ
કહેવામાં કંઈ \emph{ખોટું} નથી. આ વાત કયા આશયથી કહેવાય છે
તે માત્ર સમજવું જોઈએ. આપણે કોઈ તત્વમીમાંસાકીય તાદાત્મ્યનું
તથ્ય કહી રહ્યા નથી. આપણે ફક્ત એ નોંધીએ છીએ કે ચોક્કસ
સંદર્ભોમાં (ચોક્કસ) સંબંધોને ચોક્કસ ગણો તરીકે
\emph{લઈ શકીએ છીએ} (અને લઈશું).

\end{document}
